\documentclass[11pt,a4paper,oneside]{report}

\usepackage[a4paper,left=2.5cm,right=2.0cm,top=2.5cm,bottom=2.5cm,headheight=15pt,headsep=0.6cm,footskip=1.5cm]{geometry}
\usepackage{fontspec}
\usepackage{xeCJK}
\usepackage{setspace}
\usepackage{fancyhdr}
\usepackage{titlesec}
\usepackage{booktabs}
\usepackage{tabularx}
\usepackage{array}
\usepackage{enumitem}
\usepackage{hyperref}
\usepackage{xcolor}
\usepackage{longtable}

\setmainfont{Arial}
\setsansfont{Arial}
\setmonofont{Menlo}
\setCJKmainfont{Songti SC}
\setCJKsansfont{Hiragino Sans GB}
\setCJKmonofont{Songti SC}

\onehalfspacing
\setlength{\parindent}{2em}
\setlength{\parskip}{0.4\baselineskip}
\setlist[itemize]{leftmargin=2.5em}
\setlist[enumerate]{leftmargin=2.5em}
\hypersetup{
  colorlinks=true,
  linkcolor=black,
  citecolor=black,
  urlcolor=blue
}

\pagestyle{fancy}
\fancyhf{}
\fancyfoot[C]{\thepage}
\renewcommand{\headrulewidth}{0pt}
\fancypagestyle{plain}{
  \fancyhf{}
  \fancyfoot[C]{\thepage}
  \renewcommand{\headrulewidth}{0pt}
}

\titleformat{\chapter}[display]
  {\normalfont\bfseries\centering}
  {Chapter \thechapter.}
  {0.4em}
  {\Large}

\titleformat{name=\chapter,numberless}[display]
  {\normalfont\bfseries\centering}
  {}
  {0pt}
  {\Large}

\titleformat{\section}
  {\normalfont\bfseries}
  {\thesection.}
  {0.6em}
  {}

\titleformat{\subsection}
  {\normalfont\bfseries}
  {\thesubsection.}
  {0.6em}
  {}

\newcommand{\draftnote}[1]{\textit{#1}}

\begin{document}

\begin{titlepage}
\thispagestyle{empty}
\begin{center}
\vspace*{1.8cm}
{\Large UNDERGRADUATE SENIOR PROJECT REPORT\par}
\vspace{1.8cm}
{\Large Domain-Specific Large Language Model\\[0.4em]
Powered by Retrieval-Augmented Generation\\[0.4em]
for Scholarly Information Retrieval and Understanding\par}
\vspace{1.8cm}
by\par
\vspace{0.8cm}
Xinyu Chen (Team Leader)\\
Jian Yuan\\
Kaizhuo Hu

\vfill
Sichuan University -- Pittsburgh Institute\\
Spring 2026 Draft
\vspace{1.2cm}

Copyright in this work rests with the authors. Please ensure that any
reproduction or re-use is done in accordance with the relevant national
copyright legislation.
\end{center}
\end{titlepage}

\clearpage
\thispagestyle{empty}
\begin{center}
{\Large Declarations\par}
\end{center}
\vspace{0.8cm}

\noindent\textbf{Project Team Members}

\vspace{0.4cm}
\begin{tabularx}{\textwidth}{>{\raggedright\arraybackslash}p{3cm} >{\raggedright\arraybackslash}X >{\raggedright\arraybackslash}p{3cm} >{\raggedright\arraybackslash}p{2.5cm}}
\toprule
Role & Student Name & Student ID & Major \\
\midrule
Team Leader / Project Management & Xinyu Chen & 2022141520134 & CS \\
Frontend Development & Jian Yuan & 2022141520141 & CS \\
Database and Retrieval Module & Kaizhuo Hu & 2022141520135 & CS \\
\bottomrule
\end{tabularx}

\vspace{1.0cm}
\noindent\textbf{Advisory Committee Members}

\vspace{0.4cm}
\begin{tabularx}{\textwidth}{>{\raggedright\arraybackslash}X >{\raggedright\arraybackslash}X}
\toprule
Name & Department / Affiliation \\
\midrule
Wei Xu & Advisor \\
\draftnote{[To be completed]} & Co-Advisor / Outside Major Faculty \\
\bottomrule
\end{tabularx}

\vfill
\noindent\draftnote{This page is kept as a draft placeholder following the SCUPI template and should be revised with the final committee information before submission.}

\clearpage
\thispagestyle{empty}
\begin{center}
{\Large AI Tool Usage Statement\par}
\end{center}
\vspace{0.8cm}

\noindent This draft was prepared with AI-assisted support for language polishing,
structural organization, and preliminary drafting of literature review and system
description sections. All technical claims in the retrieval module chapter were
rechecked against the current project codebase and aligned with the team midterm
report before inclusion in this version.

\noindent\textbf{Sections involved:} Abstract, Chapter 3, and Chapter 4.

\noindent\textbf{Human post-processing:} The generated content was manually revised to
reflect the actual project status, especially the distinction between the current
retrieval-and-scoring prototype and the later full RAG pipeline planned for
subsequent development.

\noindent\draftnote{This page is a working draft only and should be replaced by the official school form before final submission.}

\clearpage
\pagenumbering{roman}
\setcounter{page}{3}

\chapter*{Abstract}
\addcontentsline{toc}{chapter}{Abstract}

This report presents a draft final thesis for the \textit{Paper Search and Analysis
System}, a scholarly search platform developed around a domain-specific paper
collection and large language model assisted interaction. The current prototype
already supports data ingestion, paper management, bilingual query preprocessing,
semantic candidate recall, and selector-based relevance scoring. In particular, the
implemented retrieval workflow centers on embedding-based paper recall with
\texttt{all-MiniLM-L6-v2}, followed by relevance filtering with the
\texttt{PASA-7B-Selector} model. Therefore, the completed core at this stage is still a
retrieval and ranking system rather than a full retrieval-augmented generation
pipeline.

This draft is intentionally focused on the retrieval module, which is also the main
scope of the author responsible for later system enhancement. To enrich the first-round
retrieval stage before the complete RAG pipeline is integrated, this report proposes a
retrieval enhancement path based on hybrid sparse-dense search, Milvus-backed vector
indexing, candidate fusion, and retrieval-oriented evaluation. The design introduces
Milvus to manage scalable vector storage and future chunk-level passage indexing,
while preserving the current selector model as the final reranking component. For
evaluation, the report prioritizes retrieval metrics such as Recall@k, MRR, nDCG,
Hit Rate, and latency; RAGAS is included as a later-stage evaluation framework once
the answer generation module is connected.

The remaining chapters follow the SCUPI senior project report structure and are kept
as concise scaffolding aligned with the midterm report. They summarize the current
project architecture, module organization, and implementation roadmap without
claiming system capabilities that have not yet been fully realized.

\noindent\textbf{Keywords:} scholarly retrieval; hybrid retrieval; Milvus; reranking;
retrieval evaluation; RAGAS

\chapter*{中文摘要}
\addcontentsline{toc}{chapter}{中文摘要}

本文给出了 \textit{Paper Search and Analysis System} 的毕业论文初稿。当前系统原型已经完成了
论文数据管理、基础检索、双语查询预处理以及基于微调模型的相关性筛选等功能。就现阶段而
言，系统主体仍然是“检索 + 排序”流程，而不是一个已经完整落地的 RAG 生成式流水线。具
体来说，当前检索主链路以 \texttt{all-MiniLM-L6-v2} 的语义向量召回为基础，再通过
\texttt{PASA-7B-Selector} 对候选论文进行打分筛选。

本稿重点完善了后续将由作者负责推进的检索增强部分。为了在完整 RAG 管线接入之前提升第
一轮召回质量，论文提前纳入了混合检索、Milvus 向量数据库、候选融合与检索评测等设计。
其中，Milvus 主要用于后续的向量索引扩展与分块级检索存储；评测方面则优先采用
Recall@k、MRR、nDCG、Hit Rate 与延迟等检索指标；RAGAS 被放在后续生成模块接入后的
联合评估阶段使用。这样既能体现系统未来的扩展方向，又不会把当前成果误写成已经完成的
完整 RAG 系统。

除检索模块外，其余章节按照中期报告和学校模板先搭建总体框架，作为后续补充实验、部署和
论文定稿的基础。

\noindent\textbf{关键词：} 学术检索；混合检索；Milvus；重排序；检索评测；RAGAS

\tableofcontents

\clearpage
\pagenumbering{arabic}

\chapter{Introduction}

\section{Background and Motivation}

The rapid growth of scholarly publications has significantly increased the burden of
literature discovery for students and researchers. Keyword-only search is often
insufficient when users formulate long, underspecified, or cross-domain research
questions. Recent progress in semantic text representation and large language models
has made it possible to move from pure lexical matching toward retrieval systems that
better capture intent, topic similarity, and citation relevance \cite{croft2010,reimers2019}.

Our project, the \textit{Paper Search and Analysis System}, is motivated by this need for a
more intelligent scholarly search workflow. The midterm stage established the basic
system infrastructure, including database construction, frontend interaction modules,
paper upload and review functions, and an initial paper retrieval engine. At the same
time, the team recognizes that the effectiveness of the overall system still depends
first on whether the right paper candidates can be found at the beginning of the
workflow.

\section{Scope of This Draft}

This draft deliberately separates the \textbf{current implementation} from the
\textbf{planned extension}. The current project status does not yet include a complete
RAG pipeline that retrieves passages and generates grounded long-form answers as the
main workflow. Instead, the implemented core is a retrieval-and-selection process:
query preprocessing, semantic candidate recall, and relevance scoring of retrieved
items by a fine-tuned selector model.

Therefore, the emphasis of this draft is not to overstate RAG as the central
achievement. Rather, the goal is to document how RAG-related retrieval ideas can be
introduced in a controlled and realistic way to strengthen the \emph{first-stage
retrieval layer}. Hybrid retrieval, Milvus-based indexing, retrieval metrics, and
RAGAS are written as the enhancement direction for the next phase, while the rest of
the report remains aligned with the real progress summarized in the midterm report.

\section{Main Contributions of This Draft}

This draft provides three concrete contributions:

\begin{enumerate}
  \item It converts the existing midterm material into a thesis-style report structure
  compatible with the SCUPI format requirements.
  \item It clarifies the real status of the current prototype by defining the retrieval
  module as the completed core and the full RAG pipeline as later work.
  \item It develops a more detailed retrieval enhancement chapter, covering hybrid
  retrieval, vector database integration with Milvus, and an evaluation plan based on
  retrieval metrics and RAGAS.
\end{enumerate}

\chapter{Related Work and Technical Background}

\section{Scholarly Retrieval and Semantic Search}

Classical information retrieval remains a strong baseline for document search, with
probabilistic ranking methods such as BM25 still widely used in practice
\cite{robertson2009}. However, semantic encoders such as Sentence-BERT make it
possible to retrieve documents by contextual similarity rather than literal term
overlap, which is especially important for academic queries containing paraphrases,
research intent, or incomplete terminology \cite{reimers2019}.

\section{Reranking and Paper Search Agents}

A practical scholarly search pipeline often benefits from multi-stage retrieval. An
initial recall stage gathers candidates quickly, and a stronger reranker is then used
to refine the result list. This project already follows that general pattern through
vector recall plus selector scoring. Recent work such as PaSa further shows that
paper-search systems can be improved through stronger decision making, selection, and
tool usage over academic corpora \cite{he2025}.

\section{Retrieval-Augmented Generation and Evaluation}

RAG combines retrieval with generation to improve factual grounding and task-specific
answer quality \cite{lewis2020}. Even so, retrieval quality remains the upstream
bottleneck. If relevant evidence is not recalled in the first place, downstream
generation cannot correct the failure. For this reason, this thesis treats retrieval
improvement as the immediate engineering priority. RAGAS is relevant here not because
the current system already completes a full RAG cycle, but because it provides a
future evaluation framework once retrieval and answer generation are connected into a
single pipeline \cite{ragas2023}.

\section{Summary}

The literature suggests a clear engineering path: combine sparse and dense retrieval,
apply stronger reranking, evaluate with explicit retrieval metrics, and only then
extend to a fully grounded generation module. The remainder of this draft follows
that logic.

\chapter{System Overview and Current Prototype}

\section{Overall Architecture}

According to the current codebase and the team midterm report, the project is
organized as a web-based paper search system with four major layers:

\begin{enumerate}
  \item a frontend interface for search, upload, library browsing, and administrator review;
  \item a Flask backend for routing, query processing, and API orchestration;
  \item a local SQLite database storing paper metadata and file references;
  \item model components for semantic encoding, selector-based scoring, and auxiliary
  LLM interaction.
\end{enumerate}

The dataset currently contains 446 papers in the local database, which is consistent
with the scale reported in the midterm report. Metadata fields include title, authors,
abstract, year, journal, DOI, review status, and citation-related fields. Optional
PDF storage is supported either locally or through Tencent Cloud COS.

\section{Functional Modules}

The current prototype is structured around the same module split described in the
midterm report:

\begin{enumerate}
  \item \textbf{Main Search Interface.} Users submit Chinese or English search requests.
  The system supports direct querying and multi-level rewritten querying.
  \item \textbf{Paper Upload Interface.} Users can upload BibTeX records and PDF files for
  later review and indexing.
  \item \textbf{Library View Interface.} Users can browse approved records and inspect
  paper metadata stored in the database.
  \item \textbf{Administrator Interface.} Administrators audit submitted entries and
  maintain data quality through review and modification operations.
\end{enumerate}

These modules already establish the engineering skeleton required for the later
retrieval enhancement and eventual full RAG integration.

\section{Current Retrieval and Scoring Workflow}

The current prototype should be described precisely. It does \textbf{not} yet run a
complete end-to-end RAG pipeline as its main workflow. Instead, its implemented
retrieval path can be summarized as follows:

\begin{enumerate}
  \item The user query is preprocessed through language detection, optional Chinese to
  English translation, spelling correction, and optional broadening or hierarchical
  rewriting.
  \item Each paper is represented mainly by title and abstract text and encoded by
  \texttt{all-MiniLM-L6-v2}.
  \item Cosine similarity is used to retrieve a candidate set from the local paper
  collection.
  \item The candidate list is sent to \texttt{PASA-7B-Selector} for relevance scoring and
  final filtering.
\end{enumerate}

This architecture means that the present system already performs semantic retrieval
and reranking, but it still operates primarily at the paper level. Although some
auxiliary PDF-based chat and citation analysis functions exist, the system has not yet
been reorganized around chunk-level evidence retrieval and grounded answer generation
as the primary interaction pattern.

\section{Current Limitations}

From the perspective of retrieval research, the main limitations of the prototype are
straightforward:

\begin{enumerate}
  \item The recall stage relies heavily on dense similarity over title and abstract
  information, which may miss documents that are lexically important but semantically
  underrepresented.
  \item The vector index is still embedded in local application logic rather than a
  dedicated vector database, which constrains scalability and future chunk-level search.
  \item The evaluation procedure has not yet been formalized into a retrieval benchmark
  with standard metrics and ablation experiments.
  \item Full RAG generation has not been integrated, so answer-level evaluation remains
  future work.
\end{enumerate}

These limitations motivate the retrieval enhancement design proposed in the next
chapter.

\chapter{Retrieval Module Enhancement Design}

\section{Design Objective}

The objective of this chapter is to improve the \textbf{first-round retrieval stage}
without mischaracterizing the project as a completed RAG system. In other words, the
purpose of the proposed enhancement is not to claim that answer generation is already
finished, but to strengthen the retrieval backbone that a later RAG pipeline will
depend on.

The design therefore keeps the current selector-based ranking logic as a useful
component and adds three new capabilities around it:

\begin{enumerate}
  \item hybrid sparse-dense retrieval for more robust candidate recall;
  \item Milvus-based vector indexing for scalable document and passage search;
  \item explicit retrieval evaluation before generation-oriented evaluation.
\end{enumerate}

\section{Planned Hybrid Retrieval Framework}

The present system mainly uses dense semantic recall. A reasonable next step is to
construct a \textbf{hybrid retrieval} architecture combining sparse lexical evidence and
dense semantic similarity. The rationale is practical: academic queries often contain
rare technical terms, benchmark names, model identifiers, and abbreviations that
lexical retrieval can preserve well, while dense retrieval is better at capturing
paraphrased intent and broader topic similarity.

The proposed hybrid retrieval workflow is shown conceptually in Table
\ref{tab:current-vs-planned}.

\begin{table}[htbp]
\centering
\caption{Current versus planned retrieval architecture}
\label{tab:current-vs-planned}
\begin{tabularx}{\textwidth}{>{\raggedright\arraybackslash}p{3.1cm} >{\raggedright\arraybackslash}X >{\raggedright\arraybackslash}X}
\toprule
Stage & Current prototype & Planned enhancement \\
\midrule
Query preprocessing & Translation, spelling correction, multi-level rewriting & Kept and further standardized for evaluation \\
Sparse retrieval & Not a dedicated stage & BM25 or equivalent lexical search over title, abstract, keywords, and chunk text \\
Dense retrieval & Sentence-transformer paper-level similarity & Dense search over paper- and chunk-level embeddings \\
Fusion & Implicitly dense-only & Weighted fusion or reciprocal rank fusion over sparse and dense candidates \\
Final ranking & PASA-7B-Selector reranking & Selector kept as the final reranker after hybrid recall \\
Generation & Not the main pipeline & Reserved for the later full RAG stage \\
\bottomrule
\end{tabularx}
\end{table}

In this design, sparse and dense retrieval will each produce a candidate list. The two
lists will then be fused before reranking. A weighted score or reciprocal rank fusion
(RRF) is a suitable strategy because it is simple, robust, and does not require large
amounts of supervised calibration data \cite{cormack2009}. This arrangement also fits
the engineering reality of the project: it improves recall quality without forcing the
team to complete the whole generation pipeline immediately.

\section{Milvus-Based Vector Index and Chunk Storage}

To support future scalability, the report proposes replacing the current in-process
vector store with a Milvus-backed index \cite{wang2021milvus}. This change is important
for two reasons.

First, Milvus provides a dedicated vector database layer for approximate nearest
neighbor search, metadata filtering, and collection management. That makes it more
suitable than ad hoc in-memory arrays when the paper collection grows beyond the
current scale. Second, Milvus offers a natural foundation for \textbf{chunk-level
retrieval}, which is necessary if the project later transitions from paper-level search
to passage-grounded response generation.

The planned indexing strategy is as follows:

\begin{enumerate}
  \item Keep a \textbf{paper-level collection} for global document recall using title,
  abstract, and structured metadata.
  \item Add a \textbf{chunk-level collection} for future passage recall, where each record
  stores paper ID, chunk ID, chunk text, section metadata, and dense embedding.
  \item Preserve metadata filters such as year, source venue, approval status, and
  domain tags to support constrained retrieval experiments.
\end{enumerate}

In the current thesis draft, chunk storage is described as a planned capability rather
than a completed result. This distinction matters. The thesis can reasonably state
that Milvus prepares the system for scalable retrieval and later RAG integration, but
it should not imply that chunk-level production deployment is already finished.

\section{Candidate Fusion and Selector-Based Reranking}

The selector model remains valuable even after hybrid recall is introduced. In fact,
the proposed enhancement keeps the selector as the \textbf{last-stage relevance
decision module}. This is consistent with the current prototype, where the selector is
already the strongest task-specific filter in the retrieval loop.

The redesigned pipeline can therefore be summarized in three stages:

\begin{enumerate}
  \item \textbf{Candidate generation.} Sparse and dense search independently retrieve a
  broad candidate pool.
  \item \textbf{Fusion.} Candidate lists are merged through RRF or a weighted combination.
  \item \textbf{Task-aware reranking.} The selector scores the fused candidates against the
  user query and returns the final list.
\end{enumerate}

This design has an important thesis-level advantage: it presents RAG-related retrieval
ideas as an \emph{incremental enhancement} to a system the team already has, rather than
as a disconnected future plan.

\section{Retrieval Metrics and Experimental Protocol}

Because the current milestone is still retrieval-centered, the primary evaluation
should also remain retrieval-centered. Table \ref{tab:retrieval-metrics} lists the
recommended metrics for the next experimental phase.

\begin{table}[htbp]
\centering
\caption{Recommended retrieval metrics for the next implementation stage}
\label{tab:retrieval-metrics}
\begin{tabularx}{\textwidth}{>{\raggedright\arraybackslash}p{3cm} >{\raggedright\arraybackslash}X}
\toprule
Metric & Purpose in this project \\
\midrule
Recall@k & Measures whether relevant papers appear in the top-k retrieved set; suitable for evaluating first-round recall improvement. \\
MRR & Captures how early the first relevant result appears; useful for interactive search experience. \\
nDCG@k & Measures ranking quality under graded relevance when some papers are more useful than others. \\
Hit Rate@k & Provides a simple success indicator for whether at least one relevant item is retrieved. \\
Latency & Measures practical usability of the retrieval pipeline after Milvus and reranking are integrated. \\
\bottomrule
\end{tabularx}
\end{table}

The benchmark construction can start from a manually curated set of representative
queries collected from the project team, advisor feedback, and the later user
acceptance testing phase. Each query should be paired with one or more relevant papers
from the approved database. Under this setup, the thesis can compare:

\begin{enumerate}
  \item dense-only retrieval;
  \item sparse-only retrieval;
  \item hybrid retrieval without selector reranking;
  \item hybrid retrieval with selector reranking.
\end{enumerate}

This evaluation plan is realistic and directly tied to the engineering work that still
needs to be completed.

\section{RAGAS in the Context of Later Full-Pipeline Evaluation}

RAGAS is included in this report because it offers a structured way to evaluate a full
RAG system without relying exclusively on expensive human annotation
\cite{ragas2023}. However, it is important to state its role carefully. Since the
current project has not yet completed the final answer-generation pipeline, RAGAS
should be treated as a \textbf{later-stage evaluation layer}, not as the primary metric
suite for the current milestone.

Once the project advances from paper retrieval to passage-grounded answer generation,
RAGAS can be used to assess aspects such as:

\begin{enumerate}
  \item context precision and context recall of the retrieved evidence;
  \item answer relevancy to the user query;
  \item faithfulness of generated responses to retrieved context.
\end{enumerate}

By describing RAGAS in this limited but forward-looking manner, the thesis keeps the
focus where it belongs: retrieval quality is the immediate problem to solve, while
generation quality will become measurable after the later RAG pipeline is actually
implemented.

\section{Relationship to the Future RAG Pipeline}

The proposed retrieval enhancement should be understood as \textbf{preparatory
infrastructure} for a future RAG system. In the later full pipeline, the retrieval
module described here will provide higher-quality paper candidates and, eventually,
chunk-level evidence passages. The downstream generation module can then consume this
evidence for explanation, summarization, citation assistance, or multi-turn paper
analysis.

Therefore, the retrieval module is not a side topic. It is the part of the project
that can be improved immediately and rigorously, while also laying the foundation for
the later answer-generation stage. This is why the present draft assigns more detail
to retrieval enhancement than to the rest of the future RAG workflow.

\chapter{Implementation Roadmap and Experimental Framework}

\section{Current Completion Status}

At the time of writing, the project has already completed the local database, the
main frontend interfaces, core upload and review functions, and the dense
retrieval-plus-selector prototype described in Chapter 3. These completed parts are
consistent with the midterm report and form the baseline for the next phase.

\section{Planned Next-Stage Tasks}

The next implementation stage can be organized into the following work packages:

\begin{enumerate}
  \item Introduce sparse retrieval over metadata and chunk text.
  \item Build Milvus collections for paper-level and later chunk-level indexing.
  \item Implement rank fusion and integrate it with the existing selector reranker.
  \item Construct a small but reliable retrieval benchmark for offline evaluation.
  \item Add the final answer-generation module after retrieval quality is stabilized.
\end{enumerate}

\section{Experimental Notes and Draft Placeholders}

\draftnote{This chapter is intentionally kept as a framework chapter in the current
draft. Final ablation tables, latency measurements, hyperparameter choices, and user
study results will be added after the retrieval enhancement is implemented and tested.}

\section{Summary}

The roadmap confirms the intended development order of the project: improve retrieval
first, evaluate retrieval explicitly, and then integrate the later full RAG pipeline
on top of a stronger retrieval backbone.

\chapter{Conclusion and Future Work}

This draft reorganizes the project into a thesis-style report and clarifies the real
status of the system. The current prototype has already established a functioning
paper search platform with semantic recall and selector-based filtering, but it should
still be described primarily as a retrieval system rather than a completed RAG
pipeline.

The main contribution of this draft is the retrieval enhancement chapter, which
proposes hybrid retrieval, Milvus-based indexing, standard retrieval metrics, and
future RAGAS-assisted evaluation. This framing is technically more accurate for the
project's current stage and also more useful for later implementation. Future work
will complete the retrieval upgrade, connect chunk-level evidence retrieval to answer
generation, and expand the report with formal experiments and deployment results.

\chapter*{References}
\addcontentsline{toc}{chapter}{References}

\begin{thebibliography}{99}

\bibitem{cormack2009}
Cormack, G. V., Clarke, C. L. A., \& B\"uttcher, S. (2009).
Reciprocal rank fusion outperforms condorcet and individual rank learning methods.
\textit{Proceedings of the 32nd International ACM SIGIR Conference on Research and Development in Information Retrieval}, 758--759.

\bibitem{croft2010}
Croft, W. B., Metzler, D., \& Strohman, T. (2010).
\textit{Search Engines: Information Retrieval in Practice}.
Addison-Wesley.

\bibitem{he2025}
He, Y., Huang, G., Feng, P., Lin, Y., Zhang, Y., Li, H., \& E, W. (2025).
PaSa: An LLM Agent for Comprehensive Academic Paper Search.
\textit{Proceedings of the 63rd Annual Meeting of the Association for Computational Linguistics}.

\bibitem{lewis2020}
Lewis, P., Perez, E., Piktus, A., Petroni, F., Karpukhin, V., Goyal, N., et al. (2020).
Retrieval-Augmented Generation for Knowledge-Intensive NLP Tasks.
\textit{Advances in Neural Information Processing Systems}, 33, 9459--9474.

\bibitem{ragas2023}
Es, S., James, J., Espinosa Anke, L., \& Schockaert, S. (2023).
RAGAS: Automated Evaluation of Retrieval Augmented Generation.
\textit{CoRR}, arXiv:2309.15217.

\bibitem{reimers2019}
Reimers, N., \& Gurevych, I. (2019).
Sentence-BERT: Sentence Embeddings using Siamese BERT-Networks.
\textit{Proceedings of the 2019 Conference on Empirical Methods in Natural Language Processing}, 3982--3992.

\bibitem{robertson2009}
Robertson, S. E., \& Zaragoza, H. (2009).
The Probabilistic Relevance Framework: BM25 and Beyond.
\textit{Foundations and Trends in Information Retrieval}, 3(4), 333--389.

\bibitem{wang2021milvus}
Wang, J., Yi, X., Guo, R., Jin, H., Xu, P., Li, S., et al. (2021).
Milvus: A Purpose-Built Vector Data Management System.
\textit{Proceedings of the 2021 International Conference on Management of Data}, 2614--2627.

\end{thebibliography}

\end{document}
